\chapter{Infantile Colic}

\section*{Definition}
Infantile colic is defined by the modified Wessel criteria as episodes of inconsolable crying and fussing in an otherwise healthy infant lasting more than 3 hours per day, more than 3 days per week, for more than 3 weeks. It typically begins around 2--3 weeks of age and resolves spontaneously by 4--6 months.

\section*{Pathophysiology}
The exact mechanism remains unclear; proposed theories include gastrointestinal immaturity (altered gut microbiome, gas production, lactose intolerance), neurodevelopmental hypersensitivity to environmental stimuli, and maternal smoking or psychosocial stress. Altered serotonin signaling and intestinal dysmotility may contribute.

\section*{Etiology}
\begin{itemize}
  \item \textbf{Gastrointestinal:} Intestinal gas, gastroesophageal reflux, cow's milk protein allergy, immature gut microbiome, lactose intolerance
  \item \textbf{Neurodevelopmental:} Overstimulation, understimulation, immature central nervous system, inability to self-soothe
  \item \textbf{Psychosocial:} Maternal anxiety, depression, smoking, inadequate caregiver responsiveness
  \item \textbf{Feeding:} Overfeeding, underfeeding, rapid feeding, improper burping, formula intolerance
  \item \textbf{Other:} Migraine equivalent, food protein-induced enterocolitis syndrome (FPIES), tobacco smoke exposure
\end{itemize}

\section*{Indications for Evaluation}
Seek care if:
\begin{itemize}
  \item Fever, lethargy, poor feeding, or inadequate weight gain (rule out serious infection or organic disease)
  \item Bilious or projectile vomiting, distended abdomen, or bloody stools
  \item Crying pattern changes or becomes higher pitched
  \item Inconsolable crying persists beyond 4--5 months of age
  \item Parental exhaustion or concern about harming the infant (risk of shaken baby syndrome)
\end{itemize}

\section*{Related Symptoms}
\begin{itemize}
  \item Crying, fussiness, clenched fists, knees drawn to chest, arching back, facial grimacing, flatus, difficulty sleeping
\end{itemize}
\textbf{Note:} Colic is a diagnosis of exclusion; red flags for serious organic causes include fever, lethargy, bilious vomiting, weight faltering, and bloody stools.

\section*{Self-Care / Home Management}
\begin{itemize}
  \item Gentle rocking, swaddling, white noise, or motion (car ride, baby swing) to provide soothing sensory input
  \item "The 5 S's": Swaddling, Side/stomach position (while awake), Shushing, Swinging, Sucking (pacifier)
  \item Burp the infant frequently during and after feeds
  \item Consider a 2-week trial of hypoallergenic formula or maternal elimination diet (cow's milk, soy, egg) if breastfed
  \item Ensure caregiver respite and stress management; never shake an infant
\end{itemize}

\section*{Diagnostic Approach}
\begin{itemize}
  \item Clinical diagnosis based on Wessel criteria; thorough history and physical exam to exclude organic causes
  \item Growth chart review to confirm adequate weight gain and normal development
  \item Urinalysis and culture to rule out UTI (especially in girls and uncircumcised boys)
  \item Stool testing for occult blood (cow's milk protein allergy) and reducing substances (lactose intolerance)
  \item Trial of acid suppression or hypoallergenic formula if GERD or milk protein allergy suspected
\end{itemize}

\section*{Treatment / Management Overview}
\begin{itemize}
  \item Reassurance and education that colic is self-limited and resolves by 4--6 months
  \item Probiotics (\textit{Lactobacillus reuteri} DSM 17938) may reduce crying time in breastfed infants
  \item Simethicone drops have not shown consistent benefit in randomized trials
  \item Avoid unproven remedies: chiropractic manipulation, herbal supplements (gripe water), and acupuncture
  \item Support caregiver mental health through validation, education, and involvement of family/friends for assistance
\end{itemize}

\clearpage
